%% ma: L12-B15-0004
%% lop: 12
%% chuong: 5
%% bai: 15
%% dang: 12.15.3
%% loai: TLN
%% muc: VD
%% dap_an: "-11"
%% nguon: "(NBV)-ĐỀ SỐ 5-MỖI NGÀY 1 ĐỀ THI 2025.pdf (D05, MỖI NGÀY 1 ĐỀ - Đề số 5), Phần III Câu 6"
%% kien_thuc: "Bài 15 (phương trình đường thẳng trong không gian, vectơ chỉ phương); Bài 8 (tích vô hướng)"
%% trang_thai: chinh-thuc
%% ghi_chu: "Giáo viên đã duyệt 10/10/2026. Đáp án gốc 3150 sai: lời giải gốc tính sai phương trình OH.u=0 (đúng là 38t-3400=0, t=1700/19). Giá trị đúng của -3a-b-c là -200/19 xấp xỉ -10,5, làm tròn là -11. Bỏ ảnh máy bay và đài kiểm soát: ảnh minh họa, không chứa dữ kiện toán"
\begin{ex}
Trong không gian $Oxyz$, đài kiểm soát không lưu sân bay có tọa độ $O(0;0;0)$, đơn vị trên mỗi trục tính theo kilômét. Một máy bay chuyển động hướng về đài kiểm soát không lưu, bay qua hai vị trí $A(-500;-250;150)$, $B(-200;-200;100)$. Khi máy bay ở gần đài kiểm soát nhất, tọa độ của vị trí máy bay là $(a;b;c)$. Giá trị của biểu thức $-3a-b-c$ là bao nhiêu (làm tròn kết quả đến hàng đơn vị)?
\shortans{-11}
\loigiai{$\overrightarrow{AB}=(300;50;-50)$ nên $\vec u=(6;1;-1)$ là một vectơ chỉ phương của đường thẳng $AB$. Máy bay gần đài nhất tại $H$ là hình chiếu vuông góc của $O$ trên $AB$: $H(-500+6t;-250+t;150-t)$.
$\overrightarrow{OH}\cdot\vec u=0\Leftrightarrow6(-500+6t)+(-250+t)-(150-t)=0\Leftrightarrow38t-3400=0\Leftrightarrow t=\dfrac{1700}{19}$.
Suy ra $H\left(\dfrac{700}{19};-\dfrac{3050}{19};\dfrac{1150}{19}\right)$, do đó $-3a-b-c=\dfrac{-2100+3050-1150}{19}=-\dfrac{200}{19}\approx-10{,}5\approx-11$.}
\end{ex}
